% Part: proof-theory
% Chapter: cut-elimination
% Section: interpolation

\documentclass[../../../include/open-logic-section]{subfiles}

\begin{document}

\olfileid[tr]{pt}{cut}{itp}

\olsection{Maehara Lemması ve Craig Aradeğerleme Teoremi}

William Craig'e ait aradeğerleme teoremi, $!A \Entails !B$ ise $!A
\Entails !C$ ve $!C \Entails !B$ olacak biçimde !!a{sentence}~$!C$ (bir
\emph{aradeğerleme}) bulunduğunu söyler. Önemli olarak,
aradeğerleme~$!C$,~$!C$ içinde oluşan !!{constant}s ile !!{predicate}s{}in
hepsi hem~$!A$ hem~$!B$ içinde oluşacak biçimde seçilebilir. (Hiçbir
!!{function}s{}ın söz konusu olmadığını varsayıyoruz.) Bu kısıtlama olmasa
$!A$~ile $!B$ önemsiz biçimde aradeğerleme görevi görebilirdi.

Aradeğerleme teoremi klasik birinci dereceden mantık için geçerlidir ve
keşfedilen ``birinci dereceden mantığın son önemli özelliği'' diye
adlandırılmıştır. Model kuramında ve otomatik akıl yürütmede önemli
uygulamaları vardır. Özgün güdüsü ve uygulaması bilim felsefesindeydi;
çünkü bir kuramın hangi ilkeller aracılığıyla aksiyomlaştırılıp
aksiyomlaştırılamayacağını incelemekte kullanılabilir. Ayrıca, bir kuram bir
kavramı örtük olarak tanımlayabiliyorsa onu açık olarak da
tanımlayabileceğini söyleyen Beth tanımlanabilirlik teoremini gerektirir.

Matematiksel mantıktaki pek çok sonuç gibi aradeğerleme teoremi de hem model
kuramının hem ispat kuramının yöntemleriyle ispatlanabilir. İspat kuramsal
yöntemlerin üstünlüğü, yalnızca bir aradeğerlemenin varlığını göstermekle
kalmayıp $!A \Entails !B$ olduğunu gösteren !!a{proof}{}tan, örneğin $!A
\Sequent !B$'nin~\Log{G3c} içindeki !!a{proof}{}ından aradeğerlemeyi hesaplayan
bir algoritma sağlamalarıdır.

$\Lang L(!A)$,~$!A$ içindeki !!{constant}s ile !!{predicate}s kümesi ve
$\Lang L(\Gamma) = \bigcup \Setabs{\Lang L(!A)}{!A \in \Gamma}$ olsun. Bu
bölüm boyunca, ele aldığımız !!{formula}s{}in hiçbir !!{function}s
içermediğini varsayıyoruz.

\begin{thm}[Craig Aradeğerleme Teoremi]\ollabel{thm:interpolation}
~$!A$'nın !!{language}{}i $\Lang L(!A)$ ve $!B$'ninki $\Lang L(!B)$ olsun.
$\Log{G3c} \Proves !A \Sequent !B$ ise, $\Log{G3c} \Proves !A \Sequent !C$
ve $\Log{G3c} \Proves !C \Sequent !B$ olacak biçimde !!a{formula}~$!C$
vardır; bu formülün !!{language}{}i $\Lang L(!C) \subseteq \Lang L(!A) \cap
\Lang L(!B)$'dır.
\end{thm}

~\Log{G3c} içindeki !!{proof}s{}ın yapısından yararlanmak için aradeğerleme
teoreminin bir genellemesini ispatlarız. Bu amaçla, sekantların öncülü ile
ardılının ayrı ayrı iki parçaya ayrıldığını düşünün. $\Gamma_1$ ile
$\Delta_1$'in birinci parçaya, $\Gamma_2$ ile $\Delta_2$'nin ikinci parçaya
ait olduğunu belirtmek için $\maeh{\Gamma_1}{\Gamma_2} \Sequent
\maeh{\Delta_1}{\Delta_2}$ yazarız. Aradeğerleme teoremi o zaman aşağıdaki
lemmadan doğrudan çıkar.

\begin{lem}[Maehara Lemması]\ollabel{lem:maehara} $\Log{G3c}
\Proves \maeh{\Gamma_1}{\Gamma_2} \Sequent \maeh{\Delta_1}{\Delta_2}$
ise, $\Log{G3c} \Proves \Gamma_1 \Sequent \Delta_1, !C$ ve $\Log{G3c}
\Proves !C, \Gamma_2 \Sequent \Delta_2$ olacak biçimde !!a{formula}~$!C$
vardır; bunun !!{language}{}i $\Lang L(!C) \subseteq \Lang L(\Gamma_1,
\Delta_1) \cap \Lang L(\Gamma_2, \Delta_2)$'dır.
\end{lem}

\begin{proof}
  $\pi \Proves \maeh{\Gamma_1}{\Gamma_2} \Sequent
  \maeh{\Delta_1}{\Delta_2}$ olduğunu varsayın. İddiayı
  $n=\pheight{\pi}$ üzerine tümevarımla ispatlarız.

  $n=0$ ise $\maeh{\Gamma_1}{\Gamma_2} \Sequent
  \maeh{\Delta_1}{\Delta_2}$ bir aksiyomdur ve atomik bir
  !!{formula}~$!A$, hem $\Gamma_1, \Gamma_2$ hem $\Delta_1, \Delta_2$ içinde
  oluşur. ~$!A$'nın solda birinci ya da ikinci parçaya ve sağda birinci ya
  da ikinci parçaya ait olmasına göre dört durum vardır.
  \begin{enumerate}
    \item $!A \in \Gamma_1$ ve $!A \in \Delta_1$. Hem $\Gamma_1 \Sequent
    \Delta_1, !C$ hem $!C, \Gamma_2 \Sequent \Delta_2$ !!{provable} olacak
    biçimde bir~$!C$ isteriz. ~$!A$ hem solda hem sağda oluştuğu için
    birincisi~$!C$'yi nasıl seçersek seçelim zaten bir aksiyomdur. $!C,
    \Gamma_2 \Sequent \Delta_2$'nın !!{provable} olmasını güvenceye almak
    için tek seçeneğimiz $!C \ident \lfalse$ almaktır.
    \item $!A \in \Gamma_1$ ve $!A \in \Delta_2$. O zaman $\Gamma_1
    \Sequent \Delta_1, !A$ ile $!A, \Gamma_2 \Sequent \Delta_2$ ikisi de
    aksiyomdur; dolayısıyla $!C \ident !A$ seçebiliriz.
    \item $!A \in \Gamma_2$ ve $!A \in \Delta_1$. O zaman $!A, \Gamma_1
    \Sequent \Delta_1$ ile $\Gamma_2 \Sequent \Delta_2, !A$ aksiyom
    olurdu; fakat~$!A$ yanlış taraflardadır. Bu nedenle $!A$ aradeğerleme
    değildir. Bununla birlikte, \RightR{\lnot} ve \LeftR{\lnot} kullanarak
    $\Gamma_1 \Sequent \Delta_1, \lnot!A$ ile $\lnot !A, \Gamma_2 \Sequent
    \Delta_2$'yı !!{prove}{}yabiliriz.
    \item $!A \in \Gamma_2$ ve $!A \in \Delta_2$. Alıştırma.
  \end{enumerate}
  $\maeh{\Gamma_1}{\Gamma_2} \Sequent \maeh{\Delta_1}{\Delta_2}$'nın
  $\lfalse \in \Gamma_1$ ya da $\lfalse \in \Gamma_2$ olduğu için bir
  aksiyom olduğu durumlar alıştırma olarak bırakılmıştır.

  $n>0$ ise $\pi$ bir mantıksal çıkarımla biter. Hangi kuralın kullanıldığına
  ve esas formülün hangi parçaya ait olduğuna göre durumları ayırmamız
  gerekir. Her durumda, lemma $\pheight{\pi_1} < n$ olan bütün
  !!{proof}s~$\pi_1$ için ve~$\pi_1$'in son-sekantı iki parçaya nasıl
  ayrılırsa ayrılsın geçerlidir biçimindeki tümevarım varsayımını
  kullanabiliriz. Bazı özel durumları ele alalım.
  \begin{enumerate}
    \item $\pi$, esas formülü~$\lnot !A$ olan $\RightR{\lnot}$ ile biter.
    İki durumumuz vardır: $\lnot !A$'nın birinci ya da ikinci parçaya ait
    olmasına göre $\pi$ şu ikisinden biriyle biter:
    \[
    \AxiomC{}
    \RightLabel{$\pi_1$}
    \Deduce$\maeh{!A, \Gamma_1}{\Gamma_2} \fCenter
      \maeh{\Delta_1}{\Delta_2}$
    \RightLabel{\RightR{\lnot}}
    \UnaryInf$\maeh{\Gamma_1}{\Gamma_2} \fCenter
      \maeh{\Delta_1, \lnot !A}{\Delta_2}$
    \DisplayProof
    \]
    ya da
    \[
    \AxiomC{}
    \RightLabel{$\pi_1$}
    \Deduce$\maeh{\Gamma_1}{!A, \Gamma_2} \fCenter
      \maeh{\Delta_1}{\Delta_2}$
    \RightLabel{\RightR{\lnot}}
    \UnaryInf$\maeh{\Gamma_1}{\Gamma_2} \fCenter
      \maeh{\Delta_1}{\Delta_2, \lnot !A}$
    \DisplayProof
    \]
    Esas !!{formula}~$\lnot !A$ birinci parçaya aitse öncüldeki etkin
    !!{formula}~$!A$'yı da birinci parçaya atadığımıza dikkat edin. Bu,
    yapabildiğimiz bir seçimdir. Etkin formülü öteki parçaya atasaydık da
    tümevarım varsayımı uygulanırdı; fakat o zaman bir aradeğerleme
    bulamazdık (alıştırma: deneyin).

    Önce $\lnot !A$'nın birinci parçaya ait olduğu ilk durumu ele alın.
    Tümevarım varsayımı gereğince~$\pi_1$'in son-sekantının bir
    aradeğerlemesi, yani !!a{formula}~$!C_1$ vardır; öyle ki aşağıdakilerin
    ikisi de
    \begin{align*}
    !A, \Gamma_1 & \Sequent \Delta_1, !C_1 &&\text{ve} &
    !C_1, \Gamma_2 & \Sequent \Delta_2
    \intertext{ispatlanabilirdir. İstediğimiz ise !!a{formula}~$!C$'dir;
    aşağıdakilerin ikisi de !!{provable} olmalıdır:}
    \Gamma_1 & \Sequent \Delta_1, \lnot !A, !C &&\text{ve} &
    !C, \Gamma_2 & \Sequent \Delta_2
    \end{align*}
    Açıkça yalnızca $!C \ident !C_1$ alabiliriz: tümevarım varsayımı sağdaki
    sekantın !!{provable} olduğunu zaten söylemiştir; soldaki sekant da
    tümevarım varsayımından aldığımız sekanta~$\RightR{\lnot}$ uygulanarak
    çıkar. Tümevarım varsayımı bize ayrıca $\Lang L(!C_1) \subseteq \Lang
    L(!A, \Gamma_1, \Delta_1) \cap \Lang L(\Gamma_2, \Delta_2)$ olduğunu
    söyler. $\Lang L(\lnot !A) = \Lang L(!A)$ ve $\Lang L(!C) = \Lang
    L(!C_1)$ olduğundan, $\Lang L(!C_1) \subseteq \Lang L(\Gamma_1,
    \Delta_1, \lnot !A) \cap \Lang L(\Gamma_2, \Delta_2)$ elde ederiz.

    İkinci durumda $\lnot !A$'nın ikinci parçaya ait olması nedeniyle durum
    biraz farklıdır. Öncüldeki etkin !!{formula}{}ü de ikinci parçaya atar ve
    tümevarım varsayımını şu sekantlara uygularız:
    \begin{align*}
    \Gamma_1 & \Sequent \Delta_1, !C_1 &&\text{ve} &
    !C_1, !A, \Gamma_2 & \Sequent \Delta_2
    \intertext{bunlar ispatlanabilirdir. Bunların ikincisine yine
    \RightR{\lnot} kuralını uygularsak şunların !!{proof}s{}ını elde ederiz:}
    \Gamma_1 & \Sequent \Delta_1, !C_1 &&\text{ve} &
    !C_1, \Gamma_2 & \Sequent \Delta_2, \lnot !A
    \end{align*}
    Bunlar, $!C_1$ bir aradeğerleme olduğunda !!{provable} olmasını
    istediğimiz sekantların tam kendisidir; yine $!C \ident !C_1$ alırız.
    Önceki gibi $\Lang L(!C_1) \subseteq \Lang L(\Gamma_1, \Delta_1) \cap
    \Lang L(\Gamma_2, \Delta_2, \lnot !A)$'dır.
    \item $\pi$, esas !!{formula}{}ü~$!A \lif !B$ olan $\RightR{\lif}$ ile
    biter. Yine $!A \lif !B$'nin birinci ya da ikinci parçaya ait olmasına
    göre iki durumumuz vardır. !!{proof}~$\pi$ şu ikisinden biriyle biter:
    \[
    \AxiomC{}
    \RightLabel{$\pi_1$}
    \Deduce$\maeh{!A, \Gamma_1}{\Gamma_2} \fCenter
      \maeh{\Delta_1, !B}{\Delta_2}$
    \RightLabel{\RightR{\lif}}
    \UnaryInf$\maeh{\Gamma_1}{\Gamma_2} \fCenter
      \maeh{\Delta_1, !A \lif !B}{\Delta_2}$
    \DisplayProof
    \]
    ya da
    \[
    \AxiomC{}
    \RightLabel{$\pi_1$}
    \Deduce$\maeh{\Gamma_1}{!A, \Gamma_2} \fCenter
      \maeh{\Delta_1}{\Delta_2, !B}$
    \RightLabel{\RightR{\lif}}
    \UnaryInf$\maeh{\Gamma_1}{\Gamma_2} \fCenter
      \maeh{\Delta_1}{\Delta_2, !A \lif !B}$
    \DisplayProof
    \]
    Etkin !!{formula}{}ü yine öncülde, sonuçtaki esas !!{formula} ile aynı
    parçaya atadık. İlk durumda tümevarım varsayımını uygular ve şunların
    !!{proof}s{}ını elde ederiz:
    \begin{align*}
    !A, \Gamma_1 & \Sequent \Delta_1, !B, !C_1 &&\text{ve} & !C_1, \Gamma_2 & \Sequent \Delta_2
    \intertext{Soldaki sekant \RightR{\lif} için uygun bir öncüldür. Böylece
    şu !!{provable} sekantlarımız vardır:}
    \Gamma_1 & \Sequent \Delta_1, !A \lif !B, !C_1 &&\text{ve} & !C_1, \Gamma_2 & \Sequent \Delta_2
    \end{align*}
    Bu,~$!C_1$'in~$\pi$'nin son-sekantı için de bir aradeğerleme olduğu
    anlamına gelir ve bir kez daha $!C \ident !C_1$ alabiliriz.
    
    Öteki durum aynı biçimde ele alınır.
    \item $\pi$, esas !!{formula}{}ü~$!A \lif !B$ olan $\LeftR{\lif}$ ile
    biter. $!A \lif !B$ birinci parçaya aitse !!{proof}~$\pi$ şöyle biter:
    \[
    \AxiomC{}
    \RightLabel{$\pi_1$}
    \Deduce$\maeh{\Gamma_1}{\Gamma_2} \fCenter
      \maeh{\Delta_1, !A}{\Delta_2}$
    \AxiomC{}
    \RightLabel{$\pi_2$}
    \Deduce$\maeh{!B, \Gamma_1}{\Gamma_2} \fCenter
      \maeh{\Delta_1}{\Delta_2}$
    \RightLabel{\LeftR{\lif}}
    \BinaryInf$\maeh{!A \lif !B, \Gamma_1}{\Gamma_2} \fCenter
      \maeh{\Delta_1}{\Delta_2}$
    \DisplayProof
    \]
    Tümevarım varsayımı şimdi bize, sol öncülün aradeğerlemesi $!C_1$ için
    iki ve ikinci öncülün aradeğerlemesi $!C_2$ için iki olmak üzere dört
    !!{provable} sekant verir:
    \begin{align*}
    \Gamma_1 & \Sequent \Delta_1, !A, !C_1 &&\text{(1)} &
    !C_1, \Gamma_2 & \Sequent \Delta_2 && \text{(2)}\\
    !B, \Gamma_1 & \Sequent \Delta_1, !C_2 &&\text{(3)} &
    !C_2, \Gamma_2 & \Sequent \Delta_2 && \text{(4)}
    \end{align*}
    Zayıflatmayı uygulayıp (1)'in sağına $!C_2$,~(3)'ün sağına $!C_1$
    eklersek, (1) ve (3) sekantları~\LeftR{\lif}'in öncülleri olabilir:
    \[
    \AxiomC{}
    \Deduce$\Gamma_1 \fCenter \Delta_1, !A, !C_1, !C_2$
    \AxiomC{}
    \Deduce$!B, \Gamma_1 \fCenter \Delta_1, !C_1, !C_2$
    \RightLabel{\LeftR{\lif}}
    \BinaryInf$!A \lif !B, \Gamma_1 \fCenter \Delta_1, !C_1, !C_2$
    \DisplayProof
    \]
    $\RightR{\lor}$ uygulayarak şu sekantı elde ederiz:
    \[
    !A \lif !B, \Gamma_1 \fCenter \Delta_1, !C_1 \lor !C_2
    \]
    Bu, $!C_1 \lor !C_2$'nin bu durumda aradeğerlememiz olabileceğini
    düşündürür. Nitekim, kalan (2) ve~(4) sekantlarından öteki gerekli
    sekantı !!{prove}{}yabiliriz:
    \[
    \AxiomC{}
    \Deduce$!C_1, \Gamma_2 \fCenter \Delta_2$
    \AxiomC{}
    \Deduce$!C_2, \Gamma_2 \fCenter \Delta_2$
    \RightLabel{\LeftR{\lor}}
    \BinaryInf$!C_1 \lor !C_2, \Gamma_2 \fCenter \Delta_2$
    \DisplayProof
    \]
    $!C_1 \lor !C_2$'nin doğru dilde olduğunu doğrulamak kalır. Tümevarım
    varsayımı bize şunları verir:
    \begin{align*}
    \Lang L(!C_1) & \subseteq
    \Lang L(\Gamma_1, \Delta_1, !A) \cap \Lang L(\Gamma_2, \Delta_2) &\text{ve}\\
    \Lang L(!C_2) & \subseteq
    \Lang L(!B, \Gamma_1, \Delta_1) \cap \Lang L(\Gamma_2, \Delta_2)
    \intertext{Şunlar nedeniyle:}
    \Lang L(!A) & \subseteq \Lang L(!A \lif !B) &\text{ve}\\
    \Lang L(!B) & \subseteq \Lang L(!A \lif !B)
    \intertext{şu ikisi de geçerlidir:}
    \Lang L(!C_1) & \subseteq
    \Lang L(\Gamma_1, \Delta_1, !A \lif !B) \cap \Lang L(\Gamma_2, \Delta_2) &\text{ve}\\
    \Lang L(!C_2) & \subseteq
    \Lang L(\Gamma_1, \Delta_1, !A \lif !B) \cap \Lang L(\Gamma_2, \Delta_2)
    \intertext{ve sonuç olarak $\Lang L (!C_1 \lor !C_2) = {}$}
    \Lang L(!C_1) \cup \Lang L(!C_2) & \subseteq
    \Lang L(\Gamma_1, \Delta_1, !A \lif !B) \cap \Lang L(\Gamma_2, \Delta_2),
    \end{align*}
    istendiği gibi.

    $!A \lif !B$'nin ikinci parçaya ait olduğu durumda durum farklıdır,
    fakat benzer biçimde ele alınır. Tümevarım varsayımı bize yine dört
    !!{provable} sekant verir:
    \begin{align*}
    \Gamma_1 & \Sequent \Delta_1, !C_1 &&\text{(1)} &
    !C_1, \Gamma_2 & \Sequent \Delta_2, !A && \text{(2)}\\
    \Gamma_1 & \Sequent \Delta_1, !C_2 &&\text{(3)} &
    !C_2, !B, \Gamma_2 & \Sequent \Delta_2 && \text{(4)}
    \end{align*}
    Şimdi (2) ve (4) sekantları---bazı zayıflatılmış !!{formula}s ile
    birlikte---~\LeftR{\lif}'in öncülleri olur:
    \[
    \AxiomC{}
    \Deduce$!C_1, !C_2, \Gamma_2 \fCenter \Delta_2, !A$
    \AxiomC{}
    \Deduce$!B, !C_1, !C_2, \Gamma_2 \fCenter \Delta_2$
    \RightLabel{\LeftR{\lif}}
    \BinaryInf$!A \lif !B, !C_1, !C_2, \Gamma_2 \fCenter \Delta_2$
    \DisplayProof
    \]
    Bunu yine $!C_1$ ile $!C_2$'den kurulmuş tek bir !!{formula} içeren bir
    sekanta dönüştürmek isteriz; fakat bu kez onu öncülde isteriz (çünkü
    artık ikinci parçayı ele alıyoruz). $\LeftR{\land}$ uygulamak işi görür;
    aradeğerleme olarak $!C \ident (!C_1 \land !C_2)$ almalıyız. Uygun bir
    rastlantıyla, öteki parça için karşılık gelen sekantı !!{prove}{}mak üzere
    kalan (1) ve (3) sekantlarını kullanabiliriz:
    \[
    \AxiomC{}
    \Deduce$\Gamma_1 \fCenter \Delta_1, !C_1$
    \AxiomC{}
    \Deduce$\Gamma_1 \fCenter \Delta_1, !C_2$
    \RightLabel{\RightR{\land}}
    \BinaryInf$\Gamma_1 \fCenter \Delta_1, !C_1 \land !C_2$
    \DisplayProof
    \]
    Önceki gibi $\Lang L(!C_1 \land !C_2) \subseteq \Lang
    L(\Gamma_1, \Delta_1) \cap \Lang L(\Gamma_2, \Delta_2, !A \lif
    !B)$'dır.

    \item $\pi$, esas !!{formula}{}ü~$\lforall[x][!A(x)]$ olan
    $\RightR{\lforall}$ ile biter:
    \[
    \AxiomC{}
    \RightLabel{$\pi_1$}
    \Deduce$\maeh{\Gamma_1}{\Gamma_2} \fCenter
    \maeh{\Delta_1, !A(c)}{\Delta_2}$
    \RightLabel{\RightR{\lforall}}
    \UnaryInf$\maeh{\Gamma_1}{\Gamma_2} \fCenter
      \maeh{\Delta_1, \lforall[x][!A(x)]}{\Delta_2}$
    \DisplayProof
    \]
    ya da
    \[
    \AxiomC{}
    \RightLabel{$\pi_1$}
    \Deduce$\maeh{\Gamma_1}{\Gamma_2} \fCenter
      \maeh{\Delta_1}{\Delta_2, !A(c)}$
    \RightLabel{\RightR{\lforall}}
    \UnaryInf$\maeh{\Gamma_1}{\Gamma_2} \fCenter
      \maeh{\Delta_1}{\Delta_2, \lforall[x][!A(x)]}$
    \DisplayProof
    \]
    İlk durumda tümevarım varsayımı bize $\Gamma_1 \Sequent \Delta_1,
    !A(c), !C_1$ ile $!C_1, \Gamma_2 \Sequent \Delta_2$'nın !!{proof}s{}ını
    verir. Birincisine $\RightR{\lforall}$ uygulayarak $\Gamma_1 \Sequent
    \Delta_1, \lforall[x][!A(x)], !C_1$'i elde edebiliriz. (Özdeğişken
    koşulu sağlanır; çünkü~$\pi$'nin sonundaki \RightR{\lforall} kuralında
    zaten sağlanmıştır.) Dolayısıyla $\Lang L(!C_1) \subseteq \Lang
    L(\Gamma_1, \Delta_1, \lforall[x][!A(x)]) \cap \Lang L(\Gamma_2,
    \Delta_2)$ koşulu sağlanıyorsa $!C_1$ bir aradeğerleme olur. Tümevarım
    varsayımı gereğince $\Lang L(!C_1) \subseteq \Lang L(\Gamma_1,
    \Delta_1, !A(c)) \cap \Lang L(\Gamma_2, \Delta_2)$'dır. O halde~$c$
    hakkında kaygılanmamız gerekir---$c$,~$!C_1$ içinde oluşabilir mi?
    Oluşamaz. $c$,~$!C_1$ içinde oluşsaydı hem $c \in \Lang L(\Gamma_1,
    \Delta_1, !A(c))$ (ki bu doğrudur) hem de $c \in \Lang L(\Gamma_2,
    \Delta_2)$ olurdu. Fakat o zaman $c$,~$\pi$ içindeki
    $\RightR{\lforall}$ çıkarımının sonucunda oluşarak özdeğişken koşulunu
    ihlal ederdi.
    
    Öteki durum benzerdir.

    \item $\pi$, esas !!{formula}{}ü~$\lexists[x][!A(x)]$ olan
    $\RightR{\lexists}$ ile biter. Yine iki durumumuz vardır.
    $\lexists[x][!A(x)]$'in birinci parçaya ait olduğu durumu ele alır,
    ötekini alıştırma olarak bırakırız.

    İlk durumda $\pi$ şöyle biter:
    \[
    \AxiomC{}
    \RightLabel{$\pi_1$}
    \Deduce$\maeh{\Gamma_1}{\Gamma_2} \fCenter
      \maeh{\Delta_1, \lexists[x][!A(x)], !A(t)}{\Delta_2}$
    \RightLabel{\RightR{\lexists}}
    \UnaryInf$\maeh{\Gamma_1}{\Gamma_2} \fCenter
      \maeh{\Delta_1, \lexists[x][!A(x)]}{\Delta_2}$
    \DisplayProof
    \]
    Tümevarım varsayımı bize $\Gamma_1 \Sequent \Delta_1,
    \lexists[x][!A(x)], !A(t), !C_1$ ile $!C_1, \Gamma_2 \Sequent
    \Delta_2$'nın !!{proof}s{}ını verir. Birincisine $\RightR{\lexists}$
    uygulayarak $\Gamma_1 \Sequent \Delta_1, \lexists[x][!A(x)], !C_1$'i
    elde edebiliriz. Dolayısıyla $\Lang L(!C_1) \subseteq \Lang
    L(\Gamma_1, \Delta_1, \lexists[x][!A(x)]) \cap \Lang L(\Gamma_2,
    \Delta_2)$ koşulu sağlanıyorsa $!C_1$ yine bir aradeğerleme olur. Burada
    tümevarım varsayımı bize $\Lang L(!C_1) \subseteq \Lang L(\Gamma_1,
    \Delta_1, \lexists[x][!A(x)], !A(t)) \cap \Lang L(\Gamma_2, \Delta_2)$
    verir. Hiçbir !!{function}s bulunmadığını varsaydığımızı şimdi
    anımsayın. Bu,~$t$ teriminin ancak !!a{constant}, yani $t \ident b$
    olabileceği anlamına gelir. $b \notin \Lang L(!C_1)$ ise ya da $b \in
    \Lang L(\Gamma_1, \Delta_1, \lexists[x][!A(x)]) \cap \Lang L(\Gamma_2,
    \Delta_2)$ ise kaygılanacak bir şey yoktur: aradeğerleme olarak $!C_1$
    alabiliriz. Peki $b \in \Lang L(!C_1)$ ve $b \notin \Lang L(\Gamma_1,
    \Delta_1, \lexists[x][!A(x)]) \cap \Lang L(\Gamma_2, \Delta_2)$ ise ne
    olur? Öncelikle,~$!C_1$ içindeki bütün~$b$ oluşumlarını
    !!a{variable}~$y$ ile değiştirerek $!C_1 \ident \Subst{!C_1'}{b}{y}$
    yazabiliriz; burada $!C_1'$~$b$ içermez. Ayrıca ya $b \notin \Lang
    L(\Gamma_1, \Delta_1, \lexists[x][!A(x)])$ ya da $b \notin \Lang
    L(\Gamma_2, \Delta_2)$'dır. Sırasıyla $!C \ident
    \lforall[y][!C_1'(y)]$ ya da $!C \ident \lexists[y][!C_1'(y)]$
    alabiliriz. Birinci durumda şunlar vardır:
    \[
    \AxiomC{}
    \Deduce$\Gamma_1 \fCenter
      \Delta_1, \lexists[x][!A(x)], !A(b), !C_1'(b)$
    \RightLabel{\RightR{\lexists}}
    \UnaryInf$\Gamma_1 \fCenter \Delta_1, \lexists[x][!A(x)], !C_1'(b)$
    \RightLabel{\RightR{\lforall}}
    \UnaryInf$\Gamma_1 \fCenter
      \Delta_1, \lexists[x][!A(x)], \lforall[y][!C_1'(y)]$
    \DisplayProof
    \]
    ve
    \[
    \AxiomC{}
    \Deduce$!C_1'(b), \lforall[y][!C_1'(y)], \Gamma_2 \fCenter \Delta_2$
    \RightLabel{\LeftR{\lforall}}
    \UnaryInf$\lforall[y][!C_1'(y)], \Gamma_2 \fCenter \Delta_2$
    \DisplayProof
    \]
    En üstteki sekantların !!{proof}s{}ı tümevarım varsayımından gelir
    (ikincisinde öncüle $\lforall[y][!C_1'(y)]$ eklemek için zayıflatmanın
    kabul edilebilirliğini uygularız). İlk !!{proof} içindeki
    \RightR{\lforall} özdeğişken koşulu, $b \notin \Lang L(\Gamma_1,
    \Delta_1, \lexists[x][!A(x)])$ olduğundan ve $!C_1'$~$b$ içermeyecek
    biçimde tanımlandığından sağlanır.
    
    İkinci durumda (tümevarım varsayımına ek olarak
    ~ $\lexists[y][!C_1'(y)]$ ile zayıflatmanın kabul edilebilirliğinden)
    şunlar vardır:
    \[
    \AxiomC{}
    \Deduce$\Gamma_1 \fCenter
      \Delta_1, \lexists[x][!A(x)], !A(b), \lexists[y][!C_1'(y)], !C_1'(b)$
    \RightLabel{\RightR{\lexists}}
    \UnaryInf$\Gamma_1 \fCenter
      \Delta_1, \lexists[x][!A(x)], \lexists[y][!C_1'(y)], !C_1'(b)$
    \RightLabel{\RightR{\lexists}}
    \UnaryInf$\Gamma_1 \fCenter
      \Delta_1, \lexists[x][!A(x)], \lexists[y][!C_1'(y)]$
    \DisplayProof
    \]
    ve
    \[
    \AxiomC{}
    \Deduce$!C_1'(b), \Gamma_2 \fCenter \Delta_2$
    \RightLabel{\LeftR{\lexists}}
    \UnaryInf$\lexists[y][!C_1'(y)], \Gamma_2 \fCenter \Delta_2$
    \DisplayProof
    \]
    İkinci !!{proof} içindeki \LeftR{\lexists} özdeğişken koşulu, $b \notin
    \Lang L(\Gamma_2, \Delta_2)$ olduğundan ve $!C_1'$~$b$ içermeyecek
    biçimde tanımlandığından sağlanır.
  \end{enumerate}
  Öteki durumlar benzerdir; onları alıştırma olarak bırakıyoruz.
\end{proof}

\begin{prob}
Şu kurallara sahip bir $\lnand$ bağlacımız olduğunu varsayın:
\[
\Axiom$\Gamma \fCenter \Delta, !A$
\Axiom$\Gamma \fCenter \Delta, !B$
\RightLabel{\LeftR{\lnand}}
\BinaryInf$!A \lnand !B, \Gamma \fCenter \Delta$
\DisplayProof
\quad
\Axiom$!A, !B, \Gamma \fCenter \Delta$
\RightLabel{\RightR{\lnand}}
\UnaryInf$\Gamma \fCenter \Delta, !A \lnand !B$
\DisplayProof
\]
~$\pi$'nin bu kurallardan biriyle bittiği \olref{lem:maehara} ispatı
durumlarını tamamlayarak Maehara Lemması'nın hâlâ geçerli olduğunu gösterin.
\end{prob}

Maehara Lemması'nı saptadıktan sonra Craig Aradeğerleme Teoremi'ni
ispatlamak için onu kullanabiliriz.

\begin{proof}[\olref{thm:interpolation} İspatı]
  $!A \Sequent !B$'nin !!{provable} olduğunu varsayın.
  Aradeğerleme~$!C$'yi elde etmek için $\maeh{!A}{\ } \Sequent
  \maeh{\ }{!B}$'ye \olref{lem:maehara}'yı uygulayın. $\Proves !A \Sequent
  !C$ ve $\Proves !C \Sequent !B$ elde ederiz.
\end{proof}

$\Gamma$, !!{sentence}s kümesi olan ve bir $n$-li yüklem~$R$ içeren,
!!a{language}~$\Lang L$ içindeki bir kuram olsun. Aşağıdaki koşul
sağlanıyorsa $\Gamma$'nın~$R$'yi \emph{örtük olarak tanımladığını} söyleriz:
\begin{align*}
\Gamma, \Gamma'
& \Entails \lforall[x_1][\dots\lforall[x_n][(R(x_1,\dots,x_n) \liff
R'(x_1, \dots, x_n))]],
\intertext{burada $\Gamma'$,~$\Gamma$ içindeki bütün~$R$ oluşumlarının
$R' \notin \Lang L$ ile değiştirildiği kuramdır. $\Gamma$,~$R$'yi ancak ve
ancak~$R$ içermeyen bir~$!A(x_1, \dots, x_n)$ var olup şu koşul
sağlandığında \emph{açık olarak tanımlar}:}
\Gamma &\Entails
\lforall[x_1][\dots\lforall[x_n][(R(x_1,\dots,x_n) \liff !A(x_1,
\dots, x_n))]].
\end{align*}
Açıkça, $\Gamma$, $R$'yi açık olarak tanımlıyorsa örtük olarak da
tanımlar. Tersi açık değildir, ancak yine de geçerlidir:

\begin{lem}[Beth Tanımlanabilirlik Teoremi]
  $\Gamma$,~$R$'yi örtük olarak tanımlıyorsa açık olarak da tanımlar.
\end{lem}

\begin{proof}
  $\Gamma$'nın~$R$'yi örtük olarak tanımladığını varsayın. $R'$, $\Gamma'$
  yukarıdaki gibi ve $\Lang L' = \Lang L(\Gamma') = \Lang L \setminus
  \{R\} \cup \{R'\}$ olsun. Varsayım gereğince şunu elde ederiz:
  \begin{align*}
    \Gamma, \Gamma' & 
    \Sequent \lforall[x_1][\dots\lforall[x_n][(R(x_1,\dots,x_n) 
    \liff R'(x_1, \dots, x_n))]]
  \intertext{$c_1$, \dots,~$c_n$,~$\Gamma$ içinde oluşmayan !!{constant}s
  olsun. Tersinirlik lemması gereğince şunu elde ederiz:}
    \Gamma, \Gamma', R(c_1,\dots,c_n) & \Sequent  R'(c_1, \dots, c_n)
  \intertext{Şuna Maehara Lemması'nı uygulayın:}
    \maeh{\Gamma, R(c_1,\dots,c_n)}{\Gamma'} & \Sequent
      \maeh{\ }{R'(c_1, \dots, c_n)}
  \intertext{Böylece aşağıdaki koşulları sağlayan bir
  ~$!A \ident !A(x_1, \dots, x_n)$ bulun:}
    \Gamma, R(c_1,\dots,c_n) & \Sequent !A(c_1, \dots, c_n) \\
    !A(c_1, \dots, c_n), \Gamma' & \Sequent R'(c_1, \dots, c_n)
  \intertext{bunların !!{proof}s{}ı vardır. $\Lang L(!A) = \Lang L_1 \cap
  \Lang L_2 = \Lang L(\Gamma) \setminus \{R\}$ olduğundan, $!A$~$R$
  içermez. İkinci sekantın !!{proof}{}ı içinde her yerde $R'$'yi~$R$ ile
  değiştirip şunun !!a{proof}{}ını elde edin:}
    !A(c_1, \dots, c_n), \Gamma & \Sequent R(c_1, \dots, c_n)
  \end{align*}
  Şunun !!a{proof}{}ını elde etmek için \RightR{\lif} ve
  \RightR{\lforall} uygulayın:
  \[
    \Gamma \Sequent \lforall[x_1][\dots\lforall[x_n][(R(x_1,\dots,x_n) \liff !A(x_1, \dots, x_n))]].
  \]
  Bu,~$!A(x_1, \dots, x_n)$'in $\Gamma$ içinde~$R$'yi açıkça
  tanımladığını gösterir.
\end{proof}

Aradeğerlemeye denk olan (ve onun gibi Maehara Lemması kullanılarak
ispatlanabilen) üçüncü bir sonuç şu teoremdir: İki tutarlı kuram~$\Gamma_1$,
$\Gamma_2$, $\Lang L(\Gamma) \subseteq \Lang L(\Gamma_1) \cap \Lang
L(\Gamma_2)$ dili içinde tam bir $\Gamma$ kuramı içeriyorsa, $\Gamma_1 \cup
\Gamma_2$ tutarlıdır.

Tam bir kuramın, kendi dilindeki her !!{sentence}~$!A$ için ya $\Gamma
\Proves !A$ ya da $\Gamma \Proves \lnot !A$ olan bir kuram olduğunu
anımsayın. $\Gamma_1$ ile~$\Gamma_2$,~$\Gamma$'nın tutarlı uzantıları
olduğundan $\Gamma$ tutarlı olmalıdır. Buradan $!A$, $\Lang L(\Gamma_1)
\cap \Lang L(\Gamma_2)$ dilinde !!a{sentence} ise $\Gamma_1 \Entails !A$
ve $\Gamma_2 \Entails \lnot !A$'nın birlikte geçerli olamayacağı çıkar.
Bunu görmek için böyle bir~$!A$ olduğunu varsayın. $\Gamma_1 \Entails !A$
ise $\Gamma \Entails !A$'dır. Aksi durumda~$\Gamma$'nın tamlığı gereğince
$\Gamma \Entails \lnot !A$ ve $\Gamma \subseteq \Gamma_1$ olduğundan
$\Gamma_1 \Entails \lnot !A$ olur; bu da $\Gamma_1$'i tutarsız kılar.
Öyleyse $\Gamma \Entails !A$ ve $\Gamma \subseteq \Gamma_2$ olduğundan
$\Gamma_2 \Entails !A$'dır. Dolayısıyla $\Gamma \Entails/ \lnot !A$;
yoksa $\Gamma_2$ tutarsız olurdu.

$\Gamma_1 \Entails !A$ ve $\Gamma_2 \Entails \lnot !A$ olan
!!a{sentence}~$!A$'nın ortak !!{language}~$\Lang L(\Gamma_1) \cap \Lang
L(\Gamma_2)$ içinde bulunmaması ispat için gereken her şeydir; şimdi ispatı
Maehara Lemması'nı kullanarak ispat kuramsal biçimde veriyoruz.

\begin{thm}[Robinson Ortak Tutarlılık Teoremi]
$\Gamma_1$, $\Gamma_2$ tutarlı kuramlar olsun ve !!{language} $\Lang
L(\Gamma_1) \cap \Lang L(\Gamma_2)$ içindeki hiçbir~$!A$ için $\Gamma_1
\Entails !A$ ve $\Gamma_2 \Entails \lnot !A$ birlikte geçerli olmasın. O
zaman $\Gamma_1 \cup \Gamma_2$ tutarlıdır.
\end{thm}

\begin{proof}
  Tersine $\Gamma_1 \cup \Gamma_2$'nin tutarsız olduğunu varsayın. O zaman
  $\Gamma_1, \Gamma_2 \Sequent \ $ sekantının !!a{proof}{}ı vardır.
  $\maeh{\Gamma_1}{\Gamma_2} \Sequent \maeh{\ }{\ }$'ye Maehara Lemması'nı
  uygulayarak, $\Proves \Gamma_1 \Sequent !A$ ve $\Proves !A, \Gamma_2
  \Sequent \quad$ olacak biçimde !!a{sentence}~$!A$ elde edin; bu tümce
  !!{language}~$\Lang L(\Gamma_1) \cap \Lang L(\Gamma_2)$ içindedir.
  İkinciden $\Proves
  \Gamma_2 \Sequent \lnot !A$'yı elde ederiz. Fakat böyle bir
  !!{sentence}{}nin bulunmadığını varsaymıştık.
\end{proof}

Yukarıda $\Gamma$, $\Gamma_1$, $\Gamma_2$ kuramlarının sonlu olduğunu
örtük olarak varsaydık. Tıkızlık gereğince sonuçlar sonsuz kuramlar için de
geçerlidir.

\end{document}
